% !TeX root = main.tex
\resetproblem
\begin{center}
    {\LARGE \textbf{Solution Manual of Problem Set 25} \par}
    \vspace{1em}
    {\large Bufan Zheng\\ \href{mailto:bufan@g.ecc.u-tokyo.ac.jp}{\texttt{bufan@g.ecc.u-tokyo.ac.jp}} \par}
\end{center}

\noindent Mostly-plus metric. Use \(T^{\mu\nu}=-\frac{2}{\sqrt{-g}}\frac{\delta S}{\delta g_{\mu\nu}}\) in curved space, evaluated at \(g=\eta\).

\begin{problem}
    \textbf{Canonical stress tensor (free scalar).} For \(\mathcal{L}=-\frac{1}{2}\partial_\mu\varphi\,\partial^\mu\varphi-\frac{m^2}{2}\varphi^2\), compute the canonical \(T^{\mu\nu}_{\text{can}}\) and verify \(\partial_\mu T^{\mu\nu}_{\text{can}}=0\) on-shell.
\end{problem}
\begin{solution}[25.1]
	The canonical momentum--energy tensor is:
	\[
	T_{\mathrm{can}}^{\mu\nu}=\frac{\partial\mathcal{L}}{\partial(\partial_\mu\varphi)}\partial^\nu\varphi-\eta^{\mu\nu}\mathcal{L}=-\partial^{\mu}\varphi\left.\partial^{\nu}\varphi-\eta^{\mu\nu}\mathcal{L}=-\partial^{\mu}\varphi\right.\partial^{\nu}\varphi+\eta^{\mu\nu}\left(\frac{1}{2}(\partial\varphi)^2+\frac{1}{2}m^2\varphi^2\right)
	\]
	The conservation law can be shown as follows:
	\[
	\partial_\mu T_{\mathrm{can}}^{\mu\nu}=-\partial_\mu(\partial^\mu\varphi\left.\partial^\nu\varphi\right)+\partial^\nu\left(\frac{1}{2}(\partial\varphi)^2+\frac{1}{2}m^2\varphi^2\right)=-(\square-m^2)\varphi\partial^\nu\varphi
	\]
	Which vanishes on-shell.
\end{solution}
\begin{problem}
    \textbf{Energy density and positivity.} Compute \(T^{00}\) for the free scalar and express it in terms of \(\dot{\varphi}\) and spatial gradients \(\nabla\varphi\). Show it is nonnegative for real \(\varphi\).
\end{problem}
\begin{solution}[25.2]
	\[
	{T^0}_0=\dot{\varphi}^2-\mathcal{L}=\frac{1}{2}\dot{\varphi}^2+\frac{1}{2}(\nabla\varphi)^2+\frac{1}{2}m^2\varphi^2\geq0
	\]
\end{solution}
\begin{problem}
    \textbf{Symmetrization via metric variation.} Minimally couple the free scalar to \(g_{\mu\nu}\). Compute \(T^{\mu\nu}\) by varying \(g_{\mu\nu}\) and show it is symmetric.
\end{problem}
\begin{solution}[25.3]
	The minimally coupled action is:
	\[
	S=\int d^Dx\sqrt{-g}\left(-\frac{1}{2}g^{\mu\nu}\nabla_\mu\varphi\nabla_\nu\varphi-\frac{1}{2}m^2\varphi^2\right)
	\]
	$$
	\begin{aligned}
		\delta S
		&=\delta\int d^Dx\,\sqrt{-g}\left(-\frac12 g^{\mu\nu}\nabla_\mu\varphi\nabla_\nu\varphi-\frac12 m^2\varphi^2\right)\\
		&=\int d^Dx\left[\delta(\sqrt{-g})\left(-\frac12 g^{\mu\nu}\nabla_\mu\varphi\nabla_\nu\varphi-\frac12 m^2\varphi^2\right)
		+\sqrt{-g}\left(-\frac12\,\delta g^{\mu\nu}\nabla_\mu\varphi\nabla_\nu\varphi\right)\right]\\
		&=\int d^Dx\,\sqrt{-g}\left[\frac12 g^{\alpha\beta}\delta g_{\alpha\beta}\left(-\frac12 g^{\mu\nu}\nabla_\mu\varphi\nabla_\nu\varphi-\frac12 m^2\varphi^2\right)
		+\frac12 g^{\mu\alpha}g^{\nu\beta}\delta g_{\alpha\beta}\nabla_\mu\varphi\nabla_\nu\varphi\right]\\
		&=\frac12\int d^Dx\,\sqrt{-g}\,\delta g_{\alpha\beta}\left[\nabla^\alpha\varphi\nabla^\beta\varphi-\frac12 g^{\alpha\beta}\left(\nabla_\rho\varphi\nabla^\rho\varphi+m^2\varphi^2\right)\right]\\
		&=\frac12\int d^Dx\,\sqrt{-g}\,\delta g_{\mu\nu}\left[\nabla^\mu\varphi\nabla^\nu\varphi-\frac12 g^{\mu\nu}\left(\nabla_\rho\varphi\nabla^\rho\varphi+m^2\varphi^2\right)\right].
	\end{aligned}
	$$
	By definition, we have:
	\[
	T^{\mu\nu}=-\nabla^\mu\varphi\operatorname{\nabla}^\nu\varphi+\frac{1}{2}g^{\mu\nu}{\left((\nabla\varphi)^2+m^2\varphi^2\right)}
	\]
	The symmetric of this quantity is obviously.
\end{solution}
\begin{problem}
    \textbf{Trace.} For \(m=0\), compute \(T^\mu{}_\mu\) in \(D\) dimensions for the minimally coupled scalar. Determine for which \(D\) it vanishes classically (without improvement).
\end{problem}
\begin{solution}[25.4]
	\[
	\left.{T^\mu}_\mu\right|_{m^2=0} = -(\nabla\varphi)^2+\frac12 \delta^\mu_\mu(\nabla\varphi)^2=\frac{D-2}{2}(\nabla\varphi)^2
	\]
	It is clear that $T^\mu_\mu$ only vanishes classically in \textbf{two dimension}.
\end{solution}
\begin{problem}
    \textbf{Improvement term.} Add the improvement \(\Delta S=\int d^Dx\,\sqrt{-g}\,\xi R\varphi^2\). In flat space, the improved stress tensor shifts by a total-derivative term. Compute the trace \(T^\mu{}_\mu\) for \(m=0\) and determine \(\xi\) such that the trace vanishes classically.
\end{problem}
\begin{solution}[25.5]
	We first calculate the variation of $\Delta S$
	\[
	\begin{aligned}\delta(\Delta S)&=\xi\int d^Dx\delta\left(\sqrt{-g}R\varphi^2\right)\\&=\xi\int d^Dx\left[\delta(\sqrt{-g})R\left.\varphi^2+\sqrt{-g}\right.\delta R\left.\varphi^2\right]\right.\\&=\xi\int d^Dx\sqrt{-g}{\left[-\frac{1}{2}g_{\mu\nu}\delta g^{\mu\nu}\right.R\left.\varphi^2+\delta R\left.\varphi^2\right]\right.}\\&=\xi\int d^Dx\sqrt{-g}{\left[-\frac{1}{2}g_{\mu\nu}R\varphi^2\delta g^{\mu\nu}+\varphi^2{\left(R_{\mu\nu}\delta g^{\mu\nu}+\left(g_{\mu\nu}\square-\nabla_\mu\nabla_\nu\right)\delta g^{\mu\nu}\right)}\right]}\\&=\xi\int d^Dx\sqrt{-g}{\left[\left(R_{\mu\nu}-\frac{1}{2}g_{\mu\nu}R\right)\varphi^2\delta g^{\mu\nu}+\varphi^2(g_{\mu\nu}\square-\nabla_\mu\nabla_\nu)\delta g^{\mu\nu}\right]}\\&=\xi\int d^Dx\sqrt{-g}{\left[G_{\mu\nu}\varphi^2\delta g^{\mu\nu}+\varphi^2g_{\mu\nu}\square\delta g^{\mu\nu}-\varphi^2\nabla_\mu\nabla_\nu\delta g^{\mu\nu}\right]}\\&\simeq\xi\int d^Dx\left.\sqrt{-g}\left[G_{\mu\nu}\varphi^2\right.\delta g^{\mu\nu}+g_{\mu\nu}\left(\square\varphi^2\right)\delta g^{\mu\nu}-\left(\nabla_\mu\nabla_\nu\varphi^2\right)\delta g^{\mu\nu}\right]\\&=\xi\int d^Dx\sqrt{-g}{\left[G_{\mu\nu}\varphi^2+g_{\mu\nu}\square\varphi^{2}-\nabla_\mu\nabla_\nu\varphi^2\right]}\delta g^{\mu\nu}.\end{aligned}
	\]
	Here, we used the following Palatini equation:
	\[
	\delta\sqrt{-g}=-\frac{1}{2}\sqrt{-g}g_{\mu\nu}\delta g^{\mu\nu},\quad
	\delta R=R_{\mu\nu}\delta g^{\mu\nu}+(g_{\mu\nu}\square-\nabla_\mu\nabla_\nu)\delta g^{\mu\nu}
	\]
	which implies:
	\[
		T_\Delta^{\mu\nu}:=-\frac{2}{\sqrt{-g}}\frac{\delta(\Delta S)}{\delta g_{\mu\nu}}=2\xi{\left[G^{\mu\nu}\varphi^2+g^{\mu\nu}\square\varphi^2-\nabla^\mu\nabla^\nu\varphi^2\right]}
	\]
	Next, we compute the trace:
	\[
	\begin{aligned}{{T_\Delta}^{\mu}}_{\mu}&=2\xi\left[g_{\mu\nu}G^{\mu\nu}\varphi^2+g_{\mu\nu}g^{\mu\nu}\square\varphi^2-g_{\mu\nu}\nabla^{\mu}\nabla^{\nu}\varphi^2\right]\\&=2\xi\left[\left(R-\frac{1}{2}g_{\mu\nu}g^{\mu\nu}R\right)\varphi^2+D\square\varphi^2-\nabla_\mu\nabla^\mu\varphi^2\right]\\&=2\xi\left[\left(1-\frac{D}{2}\right)R\varphi^2+(D-1)\square\varphi^2\right]\\&=\xi(2-D)R\varphi^2+2\xi(D-1)\square\varphi^2.\end{aligned}
	\]
	Now, come back to the flat spacetime $R=0$, we have:
	\[
	T^\mu_\mu+{{T_\Delta}^{\mu}}_{\mu}=\frac{D-2}{2}(\partial\varphi)^2+2\xi(D-1)\partial^2\varphi^{2}=\left(4\xi(D-1)+\frac{D-2}{2}\right)(\partial\varphi)^2+\cancelto{\text{OS: }0}{4\xi(D-1)\varphi\square\phi}
	\]
	so we need to choose:
	\[
	\boxed{
	\xi =-\frac{D-2}{8(D-1)}
	}
	\]
\end{solution}

\begin{problem}
    \textbf{Stress tensor conservation as Ward identity.} With a source \(J_{\mu\nu}(x)\) coupled as \(\Delta S=\frac{1}{2}\int d^Dx\,J_{\mu\nu}T^{\mu\nu}\), derive the momentum-space Ward identity for the 2-point function \(\langle TT\rangle\) that encodes \(\partial_\mu T^{\mu\nu}=0\).
\end{problem}
\begin{problem}
    \textbf{Weyl variation and tracelessness.} For an infinitesimal Weyl rescaling \(\delta g_{\mu\nu}=2\sigma g_{\mu\nu}\), show
    \[
        \delta S=-\int d^Dx\,\sqrt{-g}\,\sigma\,T^\mu{}_\mu.
    \]
    Apply this to the improved massless scalar at the conformal value of \(\xi\) and conclude \(\delta S=0\) at the classical level.
\end{problem}
\begin{solution}[25.7]
	Using the definition of energy--momentum tensor, we have:
	\[
	T^{\mu\nu}=-\frac{2}{\sqrt{-g}}\frac{\delta S}{\delta g_{\mu\nu}}\quad\Rightarrow\quad\delta S=\int d^Dx\frac{\delta S}{\delta g_{\mu\nu}}\delta g_{\mu\nu}=-\frac{1}{2}\int d^Dx\sqrt{-g}T^{\mu\nu}\delta g_{\mu\nu}
	\]
	send $\delta g_{\mu\nu}=2\sigma g_{\mu\nu}$
	\[
	\delta S=-\int d^Dx\sqrt{-g}\sigma g_{\mu\nu}T^{\mu\nu}=-\int d^Dx\sqrt{-g}\sigma T^\mu{}_\mu
	\]
	Applying the value of $\xi$ obtained in Problem 5, the classical trace ${T^{\mu}}_{\mu}$ vanishes. Consequently, the Weyl variation $\delta S_{\text{Weyl}}$ also vanishes at the classical level. In other words, there is no Weyl anomaly at the classical level.
\end{solution}
